

```latex
\section{Introduction}

Natural Language Interfaces (NLIs) have become an effective way to simplify access to complex data sources by allowing users to express information needs using natural language instead of formal query languages. At the same time, Knowledge Graphs (KGs) provide a semantic representation of information that facilitates data integration and reasoning. Combining these technologies enables users to retrieve information without requiring expertise in database systems or SPARQL.

Virtual Knowledge Graphs (VKGs) extend this idea by providing a semantic view over existing relational databases without physically transforming the data into RDF. Systems such as Ontop implement the Ontology-Based Data Access (OBDA) paradigm, allowing SPARQL queries to be automatically translated into SQL queries executed on relational databases.

This thesis focuses on adapting an existing graph-based framework to the movie domain. Starting from a relational Netflix dataset, a Virtual Knowledge Graph is constructed through ontology design and semantic mappings. The framework is then extended by investigating the integration of Large Language Models (LLMs) to support natural language interaction, user intent recognition, and recommendation capabilities.

\paragraph{Contributions}

The main contributions of this thesis are presented in Chapter~4.

\paragraph{Structure of the Thesis}

The thesis is organized as follows:

\begin{itemize}
    \item Chapter~2 introduces the theoretical background required for this work, including Knowledge Graphs, Virtual Knowledge Graphs, Ontology-Based Data Access, Ontop, RDF, SPARQL, and Large Language Models.
    \item Chapter~3 reviews the related literature on Virtual Knowledge Graphs, ontology engineering, LLM-based semantic querying, and recommendation systems.
    \item Chapter~4 presents the main contribution of this thesis, including the construction of the Virtual Knowledge Graph, ontology design, mappings, and the proposed LLM-based framework.
    \item Chapter~5 evaluates the proposed framework through experimental results and case studies.
    \item Chapter~6 concludes the thesis and discusses future research directions.
\end{itemize}
```
